% Figure: a symmetric top spinning about its own axis at rate psi-dot,
% the axis tilted at angle theta from the vertical and precessing about
% the vertical at rate phi-dot. Shows the three Euler angles and the two
% cyclic coordinates whose momenta are conserved.
\begin{figure}[htbp]
\centering
\begin{tikzpicture}[
  axis/.style={draw=enggray, thick},
  spinax/.style={draw=engblue, very thick},
  body/.style={draw=engblue, fill=engblue!12},
  arc/.style={draw=engred, -{Stealth}, thick},
  spinarc/.style={draw=engorange, -{Stealth}, thick},
  label/.style={font=\small},
]
% pivot at origin
\fill[engblack] (0,0) circle (2pt);
\node[font=\scriptsize, anchor=north east] at (0,-0.05) {pivot};

% vertical axis (space-fixed)
\draw[axis, -{Stealth}] (0,0) -- (0,5) node[anchor=south] {vertical};

% spin axis tilted by theta = 35 deg from vertical
% direction: (sin35, cos35) = (0.574, 0.819)
\coordinate (topaxis) at (2.30, 3.28); % 4 * (0.574, 0.819)
\draw[spinax] (0,0) -- (topaxis) node[anchor=south west, engblue] {spin axis};

% top body: a disc near the end of the axis, drawn as a squat ellipse
% centered at 3.2*(0.574,0.819) = (1.84, 2.62), oriented perpendicular to axis
\begin{scope}[shift={(1.84,2.62)}, rotate=35]
  \draw[body] (0,0) ellipse (0.95 and 0.30);
\end{scope}

% theta angle arc between vertical and spin axis
\draw[arc] (0,1.7) arc (90:55:1.7);
\node[engred, label] at (0.62, 1.95) {$\theta$};

% precession phi-dot: horizontal circle arc around the vertical near base
\draw[arc] (1.0,0.45) arc (0:150:1.0 and 0.32);
\node[engred, label, anchor=north] at (0.0, 0.05) {$\dot\phi$ (precession)};

% spin psi-dot: small circular arrow around the spin axis at the disc
\draw[spinarc] (2.55,3.05) arc (-30:210:0.42 and 0.16);
\node[engorange, label, anchor=west] at (2.75, 3.35) {$\dot\psi$ (spin)};

\end{tikzpicture}
\caption[Euler angles of a symmetric top]{A symmetric top pivoted at a
fixed point, described by three Euler angles: the tilt $\theta$ of the
spin axis from the vertical (red arc), the precession angle $\phi$ of
that axis about the vertical (red horizontal sweep, rate $\dot\phi$), and
the spin $\psi$ of the body about its own axis (orange, rate $\dot\psi$).
The Lagrangian depends on $\theta$ but not on $\phi$ or $\psi$, so both
$\phi$ and $\psi$ are cyclic and their conjugate momenta, the angular
momentum about the vertical and the angular momentum about the spin axis,
are conserved. Folding the two conserved momenta into an effective
potential for $\theta$ reduces the three-coordinate motion to a
one-dimensional well, whose minimum is the steady precession at constant
tilt and whose curvature sets the nutation frequency, the nodding
superposed on the sweep.}
\label{fig:vol03ch05:symmetric-top}
\end{figure}
